\documentclass[11pt]{article}

\usepackage[a4paper,margin=1in]{geometry}
\usepackage{amsmath}
\usepackage{booktabs}
\usepackage{longtable}
\usepackage{array}
\usepackage{enumitem}
\usepackage{xcolor}
\usepackage{hyperref}
\usepackage{listings}

\hypersetup{
  colorlinks=true,
  linkcolor=blue,
  urlcolor=blue
}

\lstset{
  basicstyle=\ttfamily\small,
  frame=single,
  breaklines=true,
  columns=fullflexible,
  keepspaces=true,
  showstringspaces=false
}

\newcommand{\ms}{\,\mathrm{ms}}
\newcommand{\gb}{\,\mathrm{GB}}
\newcommand{\mb}{\,\mathrm{MB}}
\newcommand{\code}[1]{\texttt{#1}}

\title{Pruning Granularity Opportunities in the Uniform Backward Pass}
\author{gpurec technical note}
\date{May 2026}

\begin{document}
\maketitle

\section{Summary}

The current backward pruning test is row-based:
\[
  a_c = \mathbf{1}\left[\max_s |R_{c,s}| \ge \tau \right],
\]
where \(R = \code{accumulated\_rhs}\) is the current adjoint
\(\partial L/\partial \Pi\), \(c\) is a clade row, \(s\) is a species, and
\(\tau\) is \code{pruning\_threshold}.  The implementation builds
\code{active\_mask} from \(\max_s |R_{c,s}|\).

The most important distinction is between the pruning criterion and the amount
of work actually skipped:

\begin{itemize}[leftmargin=*]
\item The default fused uniform path uses the row mask mainly to skip an entire
      wave when all rows in that wave are inactive.
\item If a wave has at least one active row, the default fused path generally
      still runs the fused wave kernel on the full wave.  It does compute
      \code{active\_idx} for partial waves, but that compact index is not used
      by the fused wave kernel.
\item Experimental device-pruning paths pass the mask into the Triton kernels,
      so inactive rows and inactive split parents return early, but those paths
      also launch fixed schedules for waves that the host-pruned path skipped.
      That made them slower in the 50-family benchmark.
\end{itemize}

This leads to a clear hierarchy of opportunities:

\begin{enumerate}[leftmargin=*]
\item First, make the existing default cheaper and more honest: avoid unused
      \code{nonzero} work in the fused path, and distinguish ``threshold
      inactive'' from ``work actually skipped'' in the statistics.
\item Next, test a hybrid mode: keep host whole-wave skipping, but pass
      \code{active\_mask} into fused kernels for partially active waves.
\item Then move to device-side active worklists, not fixed-schedule masked
      kernels.  Worklists preserve pruning's ability to skip work while removing
      CPU scalar decisions.
\item After row-level pruning is real, add split-level and split-side pruning.
      The post-optimization bottleneck is split-side work, so this is likely
      more valuable than further tuning the active-mask kernel itself.
\end{enumerate}

\section{Current pruning semantics}

In \code{gpurec/core/backward.py}, each backward wave \(k\) reads:
\[
  R_k = \code{accumulated\_rhs[ws:we]}.
\]
The active row mask is:

\begin{lstlisting}
clade_max = rhs.abs().max(dim=1).values
if use_pruning:
    active_mask = clade_max >= pruning_threshold
else:
    active_mask = clade_max > 0
\end{lstlisting}

With \code{use\_pruning=False}, the code still skips exact-zero rows/waves.
That is exact zero-work gating, not approximate threshold pruning.

In the normal host-pruned path, the wave loop does:

\begin{lstlisting}
active_mask = compute_active_mask(rhs_k)

if not active_mask.any():
    n_waves_skipped += 1
    n_clades_skipped += W
    continue

n_active = int(active_mask.sum().item())
n_clades_skipped += (W - n_active)

if n_active < W:
    active_idx = active_mask.nonzero(as_tuple=True)[0]
    rhs_active = rhs_k[active_idx]
\end{lstlisting}

For the non-fused fallback, \code{active\_idx} is used to compact the self-loop
solve.  For the fast fused uniform path, the call is instead:

\begin{lstlisting}
v_k = wave_backward_uniform_fused(
    Pi_star_wave, Pibar_star_wave, ws, W, S,
    dts_r, rhs_k,
    ...,
    active_mask=active_mask_for_kernels,  # None in the default host path
)
\end{lstlisting}

Since \code{active\_mask\_for\_kernels} is only set in the experimental
\code{GPUREC\_BACKWARD\_NO\_CPU\_PRUNING} or
\code{GPUREC\_DEVICE\_PRUNING} paths, partial inactive rows are generally not
masked inside the default fused kernels.

\subsection{Implication}

The default fused path has two different notions of pruning:

\begin{itemize}[leftmargin=*]
\item \textbf{Actual compute pruning}: whole inactive waves are skipped.
\item \textbf{Diagnostic row pruning}: partial inactive rows are counted in
      \code{n\_clades\_skipped}, but the fused kernels still compute them unless
      the experimental kernel-pruning flags are enabled.
\end{itemize}

This is not necessarily a correctness bug, because computing small rows is more
accurate than pruning them.  But it is a profiling and optimization trap:
\code{n\_clades\_skipped} currently overstates actual work skipped in the fused
default.

\section{Why whole-wave pruning still matters}

Whole-wave pruning is valuable because it skips the entire wave pipeline:

\begin{lstlisting}
active mask
-> DTS forward recompute, if split wave
-> fused self-loop / wave VJP
-> DTS backward accumulation, if split wave
-> uniform Pibar VJP, if split wave
\end{lstlisting}

The profiling documents show that removing pruning wholesale is very expensive.
In the 50-family wave-2 profile:

\begin{center}
\begin{tabular}{lrr}
\toprule
Variant & Median & Relative to default \\
\midrule
Current default & $157.329\ms$ & $1.00\times$ \\
Disable pruning & $241.924\ms$ & $1.54\times$ slower \\
\bottomrule
\end{tabular}
\end{center}

The no-CPU-pruning diagnostic explains the tradeoff:

\begin{center}
\begin{tabular}{lrr}
\toprule
Metric & Host pruning & Fixed-schedule device pruning \\
\midrule
\code{cudaStreamSynchronize} calls & 318 & 246 \\
D2H copies & 260 & 188 \\
Wave self-loop launches & 36 & 49 \\
Split launches & 33 & 46 \\
GPU active time & $143.400\ms$ & $149.692\ms$ \\
Outside-Nsight median & $157.329\ms$ & $161.736\ms$ \\
\bottomrule
\end{tabular}
\end{center}

So the fixed-schedule device version removed synchronization, but processed
more waves and split waves.  The extra scheduled work outweighed the sync
savings.  The conclusion is not that CPU decisions are good; it is that a
device-side design must preserve the work-skipping property of host pruning.

\section{Granularity levels}

\subsection{Wave-level pruning}

Wave-level pruning uses:
\[
  A_k = \bigvee_{c \in \mathrm{wave}(k)} a_c.
\]
If \(A_k = 0\), the whole wave is skipped.

\textbf{Strengths.}

\begin{itemize}[leftmargin=*]
\item Very high payoff per decision: one scalar decision skips multiple kernels.
\item No complicated DAG reasoning is needed.  At the time a wave is processed,
      all contributions from later waves have already reached its rows.
\item It preserves dense, contiguous wave kernels when a wave is active.
\end{itemize}

\textbf{Weaknesses.}

\begin{itemize}[leftmargin=*]
\item One active row keeps the entire wave alive.
\item It becomes less effective as more gene trees are batched together.  With
      many families, the probability that at least one row in a global wave is
      active approaches one.
\item It does not exploit partially inactive waves, which are likely common in
      large heterogeneous batches.
\end{itemize}

\textbf{Opportunity.}

For the fused default path, if we only intend to do whole-wave pruning, we do
not need \code{active\_mask.sum().item()} or \code{active\_mask.nonzero()} in
the hot path.  A cheaper whole-wave mode could compute only:
\[
  A_k = \mathbf{1}\left[\max_{c,s} |R_{c,s}| \ge \tau \right].
\]

\begin{lstlisting}
# current default fused path does more than it uses
mask = active_rows(rhs)
if not mask.any():
    continue
n_active = int(mask.sum().item())       # mostly statistics
active_idx = mask.nonzero()[0]          # unused by fused wave kernel
v_k = fused_wave(rhs, active_mask=None)

# cheaper whole-wave-only mode
wave_active = active_wave_scalar(rhs)
if not wave_active:
    continue
v_k = fused_wave(rhs, active_mask=None)
\end{lstlisting}

This would not add new pruning power, but it should reduce active-mask overhead
and remove misleading row-pruning statistics.  It is a low-risk cleanup.

\subsection{Row/clade-level pruning}

Row pruning uses the same \code{active\_mask}, but actually applies it inside
the kernels.  For row \(c\), pruning means:

\begin{itemize}[leftmargin=*]
\item skip the wave self-loop VJP for row \(c\);
\item do not accumulate parameter gradients from row \(c\);
\item skip DTS split programs whose parent is \(c\);
\item skip uniform \code{Pibar} VJP split sides whose split parent is \(c\).
\end{itemize}

This is already implemented in the experimental masked kernels.  For example,
the wave kernel checks:

\begin{lstlisting}
if USE_ACTIVE_MASK:
    row_active = tl.load(active_mask_ptr + w)
    if row_active == 0:
        store_zero_outputs()
        return
\end{lstlisting}

The DTS and \code{Pibar} kernels similarly check
\code{active\_mask[reduce\_idx[i]]} for each split parent.

\subsubsection{Most promising near-term experiment: hybrid host-wave plus row mask}

The previous failed experiment was fixed-schedule device pruning: it removed
host whole-wave decisions and launched masked kernels for all waves.  A more
promising experiment is different:

\begin{enumerate}[leftmargin=*]
\item Keep the current host whole-wave skip.
\item If the wave is active, pass \code{active\_mask} into the fused wave, DTS,
      and \code{Pibar} kernels.
\item Avoid \code{active\_idx = nonzero(mask)} unless the generic fallback needs
      compacted rows.
\end{enumerate}

\begin{lstlisting}
# proposed hybrid
mask = active_rows(rhs)
if not mask.any():
    continue

if use_fused:
    v_k = fused_wave(rhs, active_mask=mask)
    dts = dts_fused(..., active_mask=mask)
    dts_backward(..., active_mask=mask)
    pibar_vjp(..., active_mask=mask)
else:
    active_idx = mask.nonzero()[0]
    run_compacted_fallback(active_idx)
\end{lstlisting}

This preserves the benefit that made host pruning faster than fixed scheduling:
all-inactive waves are still skipped.  It also gives real row-level pruning
inside active waves.

The cost is an extra mask load and branch in the fused kernels.  That cost can
be controlled by only passing the mask when the inactive fraction is large
enough, for example:
\[
  \frac{W - n_{\mathrm{active}}}{W} \ge \rho.
\]
The threshold \(\rho\) should be benchmarked.  Reasonable first values are
\(0\), \(0.1\), \(0.25\), and \(0.5\).

\subsection{Chunk or subwave pruning}

Chunk pruning is between whole-wave and row-level pruning.  Split a large wave
into contiguous row chunks:
\[
  \mathrm{wave}(k) = B_{k,0} \cup B_{k,1} \cup \cdots
\]
and compute:
\[
  A_{k,b} = \bigvee_{c \in B_{k,b}} a_c.
\]

This can be useful because the current global wave may be active due to a
small number of rows.  A chunk-level scheduler can skip inactive chunks while
still running dense row-contiguous kernels for active chunks.

\textbf{Tradeoff.}

\begin{itemize}[leftmargin=*]
\item Smaller chunks improve pruning granularity.
\item Smaller chunks increase launch count, reduce occupancy opportunities, and
      may fragment split metadata.
\end{itemize}

The earlier \code{max\_wave\_size} sweep is relevant but not decisive:
\code{max\_wave\_size=16384} was slower than 32768 in the old 50-family run,
but that was a static memory cap, not an active chunk scheduler.  For 1000
families, chunk-level pruning may be more valuable because whole global waves
are less likely to be fully inactive.

\section{DAG subtree pruning}

It is tempting to say that if a clade is inactive, we can prune its whole
subtree.  In a tree that would often be correct.  In the clade DAG it is not
that simple.

A child clade can have multiple incoming split paths.  If one parent is
inactive, another active parent may still send adjoint to the same child.  So
we cannot delete a structural descendant subtree just because one ancestor row
is inactive.

The safe dynamic rule is:

\begin{quote}
Process waves in reverse topological order.  When a row's wave is reached,
all already-processed parents have accumulated their contributions into that
row's RHS.  If the final row RHS is below threshold, that row can be pruned.
\end{quote}

This is exactly what row-level RHS pruning approximates.  A more aggressive
``subtree'' implementation should therefore be phrased as an active-dependency
worklist, not as a static deletion of descendants:

\begin{lstlisting}
active_rows_k = rows in wave k with absmax(rhs[row]) >= threshold
active_splits_k = splits whose parent row is in active_rows_k

for each active split:
    accumulate into child RHS

# child rows become active later only if their accumulated RHS is large enough
\end{lstlisting}

In other words, the DAG pruning object is not a subtree.  It is the dynamically
reachable set of rows under the current adjoint threshold.

\section{Split-level pruning}

Row-level pruning skips all splits of inactive parents.  It does not skip
individual negligible splits of an active parent.  That matters because the
post-optimization profile is again dominated by split-side work:

\begin{center}
\begin{tabular}{lrr}
\toprule
Component after proposal 6 & Approximate Nsys bucket & Notes \\
\midrule
\code{\_dts\_cross\_backward\_accum\_kernel} & $32.606\ms$ & split-major direct adjoints \\
\code{\_uniform\_cross\_pibar\_vjp\_tree\_from\_ud\_kernel} & $25.618\ms$ & two split sides \\
\code{\_dts\_fused\_kernel} & $12.207\ms$ & DTS forward recompute \\
\bottomrule
\end{tabular}
\end{center}

The largest split waves are root-like:

\begin{center}
\begin{tabular}{lrrr}
\toprule
Wave shape & Rows & Splits & Fanout \\
\midrule
50-family wave example & 247 & 42155 & 170.7 splits/parent \\
50-family wave example & 39 & 24229 & 621.3 splits/parent \\
\bottomrule
\end{tabular}
\end{center}

If a parent row is active, row-level pruning keeps all its splits.  For these
root-like waves, split-level pruning could be much more selective.

\subsection{Split predicates}

A split \(i\) with parent \(p(i)\) should at least satisfy:
\[
  b_i = a_{p(i)}.
\]
That is parent-row pruning.

A stronger approximate split predicate would test the actual split
responsibility or VJP contribution:
\[
  b_i =
  a_{p(i)}
  \wedge
  \left[
    \max_s |v_{p(i),s}| \cdot \omega_{i,s} \ge \tau_{\mathrm{split}}
  \right],
\]
where \(\omega_{i,s}\) is the split's normalized contribution to the parent DTS
logsumexp for species \(s\).

The challenge is that \(\omega_{i,s}\) is not free.  It depends on the same
DTS terms that the split kernels compute.  Therefore split pruning should be
introduced in stages:

\begin{enumerate}[leftmargin=*]
\item First use parent-active split pruning, because that only needs
      \code{active\_mask[reduce\_idx]}.
\item Then collect statistics on per-split contribution magnitudes inside the
      existing DTS kernel, without changing results.
\item Only if many active-parent splits have tiny contribution, add a compacted
      active-split worklist.
\end{enumerate}

\subsection{Split-side pruning}

The staged \code{Pibar} path computes the tree input:
\[
  u_d = \mathrm{grad\_Pibar} \cdot \mathrm{inv\_denom},
  \qquad
  A = \sum_s u_d[s].
\]
This creates a natural split-side pruning criterion:
\[
  d_{i,\mathrm{side}} =
  \mathbf{1}\left[
    \max_s |u_{d,i,\mathrm{side},s}| \ge \tau_{\mathrm{side}}
  \right].
\]

If \(d_{i,\mathrm{side}} = 0\), the corresponding \code{Pibar} tree VJP side
can be skipped.  This is attractive because the current staged tree VJP bucket
is about $25.6\ms$ in the post-proposal-6 profile.

The staged DTS kernel already touches the values needed to compute this
predicate.  The missing piece is scheduling: the tree kernel should run only on
active split sides rather than launching a fixed \(2 n_{\mathrm{ws}}\) grid.

\section{Species-lane or tile pruning}

Species-lane pruning would skip individual species entries or species tiles
inside an active row.  This is the finest granularity:
\[
  a_{c,tile} =
  \mathbf{1}\left[
    \max_{s \in tile} |R_{c,s}| \ge \tau_{\mathrm{tile}}
  \right].
\]

This is currently the least attractive pruning level.

\begin{itemize}[leftmargin=*]
\item Uniform \code{Pibar} uses row reductions and species-tree corrections.
      Dropping lanes changes row sums and ancestor/subtree corrections.
\item The wave kernel's expensive work is dense in \(S\), and many formulas use
      \(\sum_s\) or max over the full row.
\item A row can have many species lanes just below a max threshold.  Their
      aggregate contribution can still matter for parameter gradients.
\item The active-mask kernel already streams all \(S\) lanes once; additional
      tile metadata may not pay for itself unless rows are extremely sparse in
      species space.
\end{itemize}

Species pruning should therefore be considered only after row and split
worklists are measured.  A safer intermediate is to improve the row predicate,
for example by recording both \(\|R_c\|_\infty\) and \(\|R_c\|_1\), rather than
dropping species lanes.

\section{Family-level and batched-tree granularity}

The current batching strategy builds large global waves across many gene trees.
That improves occupancy, but it weakens whole-wave pruning.  With many families,
a global wave is active if any family contributes an active row:
\[
  A_k^{global} =
  \bigvee_g \bigvee_{c \in \mathrm{wave}(k,g)} a_c.
\]

As the number of families grows, \(A_k^{global}\) tends to be true even if most
families are inactive at that depth.  This is likely important for 1000-tree
batches.

The pruning-aware alternative is to schedule at a family-wave or chunk level:

\begin{lstlisting}
for global wave k:
    build active row counts by family or row chunk
    select active family-wave chunks
    launch kernels only for selected chunks
\end{lstlisting}

This creates a three-way tradeoff:

\begin{itemize}[leftmargin=*]
\item Larger batches improve occupancy and reduce launch overhead.
\item Smaller family/chunk units improve pruning selectivity.
\item Buffer lifetime and memory pressure constrain how many chunks can be
      resident at once.
\end{itemize}

For 1000 trees, this may matter more than further optimizing the row mask
kernel.  The current active-mask kernel is already a simple bandwidth-streaming
kernel; the larger question is whether its output can drive a better scheduler.

\section{Recommended optimization sequence}

\subsection{P0: fix measurement and remove unused default work}

Before changing pruning behavior, separate statistics into:

\begin{itemize}[leftmargin=*]
\item rows below threshold;
\item rows whose self-loop work was actually skipped;
\item split programs skipped;
\item \code{Pibar} split sides skipped;
\item whole waves skipped.
\end{itemize}

Then remove \code{active\_idx = active\_mask.nonzero()} from the fused default
unless a compact fallback path will actually use it.  If production does not
need \code{n\_clades\_skipped} every iteration, also avoid
\code{active\_mask.sum().item()} in the hot path or compute it asynchronously
for diagnostics.

\subsection{P1: hybrid host-wave skip plus fused row masks}

This is the highest-value low-complexity experiment:

\begin{lstlisting}
mask = active_rows(rhs)
if not mask.any():
    skip whole wave
else:
    pass mask into fused wave, DTS forward, DTS backward, and Pibar VJP kernels
\end{lstlisting}

It should be tested with several policies:

\begin{center}
\begin{tabular}{ll}
\toprule
Policy & Meaning \\
\midrule
always mask active waves & Maximize row pruning, measure branch overhead \\
mask if \(n_{\mathrm{active}} < W\) & Avoid all-active mask overhead \\
mask if inactive fraction \(\ge 10\%\) & Avoid overhead on nearly full waves \\
mask only split waves & Target DTS and \code{Pibar} buckets first \\
\bottomrule
\end{tabular}
\end{center}

This experiment is different from the failed device-pruning prototype because
it keeps whole-wave skipping.

\subsection{P2: device active row and split worklists}

If P1 shows meaningful inactive-row work inside active waves, the next step is
a worklist scheduler:

\begin{lstlisting}
row_mask = active_rows(rhs)
active_rows = select(rows, row_mask)
active_splits = select(splits, row_mask[reduce_idx])
active_sides = select(split_sides, row_mask[reduce_idx] & side_active)

wave_kernel_active(active_rows)
dts_kernel_active(active_splits)
pibar_tree_active(active_sides)
\end{lstlisting}

This can be built with CUB-style prefix sums/select operations or a C++ CUDA
extension.  The key is to avoid the fixed-schedule failure mode: inactive waves
and inactive split groups must not merely early-return after launch; they
should be absent from the launched work.

\subsection{P3: split-side pruning after instrumentation}

Instrument the staged DTS path to record:

\begin{itemize}[leftmargin=*]
\item max and sum of \(|u_d|\) per split side;
\item max VJP contribution per split;
\item active split count under candidate thresholds;
\item how these counts correlate with final gradient error.
\end{itemize}

If many split sides are negligible, compact \code{Pibar} tree VJP to active
sides.  This directly targets the roughly $25.6\ms$ staged tree bucket in the
post-proposal-6 profile.

\subsection{P4: family/chunk-aware scheduling for large batches}

For 1000 trees, record active rows by family and wave.  The relevant metric is
not only the number of active global waves, but:
\[
  \frac{\text{active rows in wave}}{\text{rows scheduled in wave}}.
\]

If global waves are almost always active but sparse by family or chunk, use
family-wave chunks or row chunks as the scheduling unit.  This may recover
pruning selectivity without sacrificing all batching benefits.

\section{Expected gains and risks}

\begin{longtable}{p{0.20\linewidth}p{0.22\linewidth}p{0.22\linewidth}p{0.26\linewidth}}
\toprule
Opportunity & Expected gain & Risk & Why \\
\midrule
\endfirsthead
\toprule
Opportunity & Expected gain & Risk & Why \\
\midrule
\endhead

Remove unused \code{nonzero}/stats in fused default & Small, likely $1$--$3\ms$ if partial waves are common & Low & Does not change numerical pruning semantics if statistics are deferred or corrected. \\

Host-wave skip plus active mask in fused kernels & Medium, data-dependent & Medium & Preserves whole-wave skip while enabling real row/split-parent pruning; changes approximate pruning results for rows below threshold. \\

Device active worklists & Medium to high, especially for 1000 trees & High & Removes CPU syncs and no-op work, but requires robust compaction, buffer management, and irregular-kernel interfaces. \\

Split-side \code{Pibar} pruning & Potentially high if many sides are tiny & Medium/high & Targets a $25\ms$-scale bucket, but needs new contribution thresholds and active-side scheduling. \\

Family/chunk scheduling & Potentially high for large batches & High & Whole global waves become less selective as family count grows; chunking restores selectivity but complicates memory/lifetime scheduling. \\

Species-tile pruning & Low/uncertain & High & Uniform \code{Pibar} and parameter reductions are dense over species; correctness bounds are harder. \\

\bottomrule
\end{longtable}

\section{Bottom line}

The current default mainly prunes whole waves.  That is still valuable, but it
leaves three forms of work on the table:

\begin{enumerate}[leftmargin=*]
\item partially inactive rows inside active waves;
\item inactive or negligible splits of active split waves;
\item inactive family/chunk regions inside large cross-family waves.
\end{enumerate}

The failed device-pruning experiment should not discourage pruning work.  It
tested the wrong tradeoff: it removed host decisions but gave up whole-wave
skipping.  The next pruning experiments should preserve whole-wave skipping
while adding finer-grained masks, then graduate to device worklists if the
active fractions justify the complexity.

\end{document}
